\chapter{Longueurs, angles et lieux géométriques}
\section*{Exercice 17.1 -- Al-Kashi}
Al-Kashi donne $a^2=4^2+6^2-2\cdot4\cdot6\cos120^\circ=16+36+24=76$. Donc $a=2\sqrt{19}$. L'angle obtus entraîne ici un carré supérieur à la somme des deux autres carrés.

\section*{Exercice 17.2 -- Milieu et lieu}
L'identité du milieu donne $MI^2=k+9$. Pour $k=0$, c'est le cercle de centre $I$ et rayon $3$. Pour $k=-9$, c'est le seul point $I$. Pour $k=-10$, aucune solution réelle puisque le carré serait négatif.

\section*{Exercice 17.3 -- Médiane}
Si $I$ est le milieu de $[BC]$, l'identité donne $AB^2+AC^2=2AI^2+BC^2/2$. Ainsi $74=2AI^2+32$, donc $AI^2=21$ et $AI=\sqrt{21}$.

\section*{Exercice 17.4 -- Aire par un angle}
La hauteur relative au côté de longueur $5$ vaut $8\sin(\pi/6)=4$. L'aire est donc $5\cdot4/2=10$. C'est la formule $\aire=bc\sin A/2$, obtenue en projetant un côté sur la direction perpendiculaire à l'autre.

\section*{Exercice 17.5 -- Centre de gravité}
La moyenne des coordonnées donne $G=(2,1)$. La somme des vecteurs $\overrightarrow{GA},\overrightarrow{GB},\overrightarrow{GC}$ est nulle. En développant chaque distance depuis $G$, la somme vaut $3MG^2+GA^2+GB^2+GC^2$, minimale uniquement pour $M=G$.

\section*{Exercice 17.6 -- Hauteurs sans dessin}
La hauteur issue de $C$ est verticale : $x=1$. Le côté $AC$ a pour vecteur $(1,3)$ ; la hauteur issue de $B$ vérifie $(x-4)+3y=0$. À $x=1$, on trouve $y=1$, donc $H=(1,1)$. Enfin $\overrightarrow{AH}=(1,1)$ et $\overrightarrow{BC}=(-3,3)$ ont un produit nul : la troisième hauteur passe bien par $H$.
